\subsection{A pairing obstruction to enlarging the rational torsion subgroup}
\label{sec:lang-pairing-rank-bound}

\paragraph{Finding A25 (a prime-level Lang test forces torsion rank at most two).}
Let $C/\mathbf F_Q$ be a genus-two curve in characteristic two,
let $J$ be its Jacobian, and put $G=J(\mathbf F_Q)$.
Let $r$ be an odd prime and
$\kappa=\dim_{\mathbf F_r}G[r]$.
Suppose either the nonzero-point numerical test (LW.2), or
the weaker ordinary-case positivity test following (LW.1),
holds for the subgroup $H$ under consideration.
Then the following implications hold:
\[
\begin{array}{c|c|c}
\text{hypothesis on }2\pmod r&H&\text{range forcing }\kappa\le2\\
\hline
\operatorname{ord}_r(2)=r-1&G[r]&r\ge11\\
\operatorname{ord}_r(2)=r-1&G[3r]&r\ge17\\
\langle2,-1\rangle=(\mathbf Z/r)^\times&G[r]&r\ge29\\
\langle2,-1\rangle=(\mathbf Z/r)^\times&G[3r]&r\ge41.
\end{array}
\tag{LP.1}
\]
In each range, writing $s=\sqrt Q$ and $J_Q=|G|$, the full
test (LW.2) therefore requires respectively
\[
 r>\left(J_Q\frac{2s+1}{(s-1)^2}\right)^{1/2}
 \quad\text{or}\quad
 r>\frac19\left(J_Q\frac{2s+1}{(s-1)^2}\right)^{1/2}.
\tag{LP.2}
\]
These thresholds have asymptotic sizes
$\sqrt2 Q^{3/4}$ and $(\sqrt2/9)Q^{3/4}$.
The corresponding ordinary thresholds replace $2s+1$ by $2s$.
These are necessary conditions for the specified sufficient
numerical tests, not for the existence of theta points.

\begin{proof}
The canonical principal polarization of the Jacobian gives a
nondegenerate alternating pairing
\[
                e_r:J[r]\times J[r]\longrightarrow\mu_r
\tag{LP.3}
\]
on its four-dimensional geometric $r$-torsion space, equivariant
for Frobenius.  For the pairing and its equivariance see
\cite[\S16, p.~29, and Lemma~16.2(e), p.~30]{MilneAV};
the principal polarization identifies the dual torsion space
with the original one.
If $\kappa\ge3$, the subspace $G[r]$ is not totally isotropic:
a totally isotropic subspace $W$ satisfies
$W\subseteq W^\perp$ and hence
$\dim W\le4-\dim W$.
Choose $P,T\in G[r]$ with $e_r(P,T)\ne1$.
Since $r$ is prime, this pairing value has exact order $r$.
Both points are fixed by Frobenius, so equivariance gives
\[
                         e_r(P,T)^Q=e_r(P,T).
\]
It follows that $r\mid Q-1$.
Writing $Q=2^v$ and $d=\operatorname{ord}_r(2)$, we obtain
\[
                         d\mid v,\qquad Q\ge2^d.
\tag{LP.4}
\]

The weaker of the two Lang tests in the statement already
implies
\[
 |H|>\frac{2sJ_Q}{(s-1)^2}
       \ge 2s(s-1)^2,
\tag{LP.5}
\]
where the last inequality uses $J_Q\ge(s-1)^4$.
For $Q\ge16$, the last expression is at least
$(9/8)Q^{3/2}$, because
$2(1-1/s)^2\ge2(3/4)^2=9/8$.
In particular, either successful test requires
$|H|>Q^{3/2}$ in this range.

If $H=G[r]$, then $|H|=r^\kappa\le r^4$.
If $H=G[3r]$, the ranges in (LP.1) exclude $r=3$, and
primary decomposition gives
$|H|=|G[3]|\,|G[r]|\le81r^4$.
If $\kappa\ge3$, (LP.4) now forces failure of (LP.5)
whenever
\[
                         c r^4<2^{3d/2},
\tag{LP.6}
\]
where $c=1$ or $81$ respectively.

For a primitive root, $d=r-1$.
The two starting inequalities are
\[
 11^4=14641<2^{15},\qquad
 81\cdot17^4=6765201<2^{24}.
\]
They persist on increasing the odd integer $r$ by two:
the left side grows by a factor at most
$(13/11)^4<2$, whereas the right side grows by $8$.
Thus (LP.6) holds in the first two ranges of (LP.1).

Under the signed-doubling hypothesis,
$r-1=|\langle2,-1\rangle|\le2d$, so
$d\ge(r-1)/2$.
For the other two ranges it suffices to check
\[
 29^4=707281<2^{21},\qquad
 81\cdot41^4=228886641<2^{30}.
\]
On increasing $r$ by two, the lower bound
$2^{3(r-1)/4}$ on the right of (LP.6) grows by
$2^{3/2}>2$, whereas the left side again grows by a
factor less than $2$.
All four ranges also imply $Q\ge16$ through (LP.4).
Consequently $\kappa\ge3$ contradicts either successful
Lang test, proving (LP.1).

Finally, solving (LW.2) for the subgroup size gives
$|H|>J_Q(2s+1)/(s-1)^2$.
With $\kappa\le2$, one has $|G[r]|\le r^2$ and
$|G[3r]|\le81r^2$, which proves (LP.2).
The asymptotics follow from $J_Q=Q^2+O(Q^{3/2})$.
\end{proof}

\paragraph{Implication for the proposed arithmetic route.}
In particular, first passing to an extension where all of
$J[r]$ is rational cannot force the prime-level Lang test
for an Artin prime $r\ge11$: it creates $\kappa=4$, the
excluded situation in (LP.1).
The factor $|G[3]|\le81$ can improve an individual index,
but the analogous strategy with $G[3r]$ fails for Artin
primes $r\ge17$.
This argument concerns these rational subgroup estimates.
It does not remove geometric theta witnesses, and does not
assert the same rank bound for truncated prime-power
subgroups $G[r^b]$ without an additional argument.
